% GENERATED FILE. Edit canonical lessons, then run npm run generate:books.
% canonical-source-hash: fnv1a64:bd0b2a7a5d50b949
% canonical-lessons: TA-C168-ulu, TA-C168-tontu, TA-C168-nittu, TA-C168-valukku, TA-C168-kottavi-vitu

\chapter{To plough, To dig, To stretch out, To slip, To yawn}
\label{ch:ta-to-plough-to-dig-to-stretch-out-to-slip-to-yawn}
\hlchaptermodality{168}

\begin{chapteropening}
\textbf{By the end of this chapter:} \emph{I can say to plough, to dig, to stretch out, to slip and to yawn.}

Complete the last lesson of chapter 168: I can say to plough, to dig, to stretch out, to slip and to yawn.
\end{chapteropening}

\section[\texorpdfstring{\ta{உழு} (uḻu)}{உழு (uḻu)}]{\ta{உழு} (uḻu) — to plough}

\label{lesson:TA-C168-ulu}



\begin{quote}
Before the new one: say the Tamil for to hang, then the Tamil for to roll.
\end{quote}

\subsection*{You'll want to know: \ta{உழு}}
\textbf{\ta{உழு}} — \emph{uḻu} — \textquotedblleft{}to plough\textquotedblright{}.

\begin{grammarlens}[title={What you've built}]
One more word for this chapter.
\end{grammarlens}

\subsection*{Your turn}
\begin{itemize}
\raggedright
  \item \emph{Say it:} \emph{uḻu}
  \item \emph{Say it:} \emph{uḻu}, once more
  \item \emph{Say it:} say \emph{uḻu}
  \item \emph{Recall:} say the Tamil for to hang, then the Tamil for to roll, then say \emph{uḻu} again
\end{itemize}

\begin{tcolorbox}[breakable,colback=teal!4,colframe=teal!35!black,title={Before you move on}]
What does \ta{உழு} mean? (\textquotedblleft{}To plough\textquotedblright{}.) Say it once more.
\end{tcolorbox}
\section[\texorpdfstring{\ta{தோண்டு} (tōṇṭu)}{தோண்டு (tōṇṭu)}]{\ta{தோண்டு} (tōṇṭu) — to dig}

\label{lesson:TA-C168-tontu}



\begin{quote}
Before the new one: say the Tamil for to roll, then the Tamil for to plough.
\end{quote}

\subsection*{You'll want to know: \ta{தோண்டு}}
\textbf{\ta{தோண்டு}} — \emph{tōṇṭu} — \textquotedblleft{}to dig\textquotedblright{}.

\begin{grammarlens}[title={What you've built}]
One more word for this chapter.
\end{grammarlens}

\subsection*{Your turn}
\begin{itemize}
\raggedright
  \item \emph{Say it:} \emph{tōṇṭu}
  \item \emph{Say it:} \emph{tōṇṭu}, once more
  \item \emph{Say it:} say \emph{tōṇṭu}
  \item \emph{Recall:} say the Tamil for to roll, then the Tamil for to plough, then say \emph{tōṇṭu} again
\end{itemize}

\begin{tcolorbox}[breakable,colback=teal!4,colframe=teal!35!black,title={Before you move on}]
What does \ta{தோண்டு} mean? (\textquotedblleft{}To dig\textquotedblright{}.) Say it once more.
\end{tcolorbox}
\section[\texorpdfstring{\ta{நீட்டு} (nīṭṭu)}{நீட்டு (nīṭṭu)}]{\ta{நீட்டு} (nīṭṭu) — to stretch out, to extend}

\label{lesson:TA-C168-nittu}



\begin{quote}
Before the new one: say the Tamil for to plough, then the Tamil for to dig.
\end{quote}

\subsection*{You'll want to know: \ta{நீட்டு}}
\textbf{\ta{நீட்டு}} — \emph{nīṭṭu} — \textquotedblleft{}to stretch out, to extend\textquotedblright{}.

\begin{grammarlens}[title={What you've built}]
One more word for this chapter.
\end{grammarlens}

\subsection*{Your turn}
\begin{itemize}
\raggedright
  \item \emph{Say it:} \emph{nīṭṭu}
  \item \emph{Say it:} \emph{nīṭṭu}, once more
  \item \emph{Say it:} say \emph{nīṭṭu}
  \item \emph{Recall:} say the Tamil for to plough, then the Tamil for to dig, then say \emph{nīṭṭu} again
\end{itemize}

\begin{tcolorbox}[breakable,colback=teal!4,colframe=teal!35!black,title={Before you move on}]
What does \ta{நீட்டு} mean? (\textquotedblleft{}To stretch out, to extend\textquotedblright{}.) Say it once more.
\end{tcolorbox}
\section[\texorpdfstring{\ta{வழுக்கு} (vaḻukku)}{வழுக்கு (vaḻukku)}]{\ta{வழுக்கு} (vaḻukku) — to slip}

\label{lesson:TA-C168-valukku}



\begin{quote}
Before the new one: say the Tamil for to dig, then the Tamil for to stretch out.
\end{quote}

\subsection*{You'll want to know: \ta{வழுக்கு}}
\textbf{\ta{வழுக்கு}} — \emph{vaḻukku} — \textquotedblleft{}to slip\textquotedblright{}.

\begin{grammarlens}[title={What you've built}]
One more word for this chapter.
\end{grammarlens}

\subsection*{Your turn}
\begin{itemize}
\raggedright
  \item \emph{Say it:} \emph{vaḻukku}
  \item \emph{Say it:} \emph{vaḻukku}, once more
  \item \emph{Say it:} say \emph{vaḻukku}
  \item \emph{Recall:} say the Tamil for to dig, then the Tamil for to stretch out, then say \emph{vaḻukku} again
\end{itemize}

\begin{tcolorbox}[breakable,colback=teal!4,colframe=teal!35!black,title={Before you move on}]
What does \ta{வழுக்கு} mean? (\textquotedblleft{}To slip\textquotedblright{}.) Say it once more.
\end{tcolorbox}
\section[\texorpdfstring{\ta{கொட்டாவி} \ta{விடு} (koṭṭāvi viṭu)}{கொட்டாவி விடு (koṭṭāvi viṭu)}]{\ta{கொட்டாவி} \ta{விடு} (koṭṭāvi viṭu) — to yawn}

\label{lesson:TA-C168-kottavi-vitu}



\begin{quote}
Before the new one: say the Tamil for to stretch out, then the Tamil for to slip.
\end{quote}

\subsection*{You'll want to know: \ta{கொட்டாவி} \ta{விடு}}
\textbf{\ta{கொட்டாவி} \ta{விடு}} — \emph{koṭṭāvi viṭu} — \textquotedblleft{}to yawn\textquotedblright{}.

\begin{grammarlens}[title={What you've built}]
One more word for this chapter.
\end{grammarlens}

\subsection*{Your turn}
\begin{itemize}
\raggedright
  \item \emph{Say it:} \emph{koṭṭāvi viṭu}
  \item \emph{Say it:} \emph{koṭṭāvi viṭu}, once more
  \item \emph{Say it:} say \emph{koṭṭāvi viṭu}
  \item \emph{Recall:} say the Tamil for to stretch out, then the Tamil for to slip, then say \emph{koṭṭāvi viṭu} again
\end{itemize}

\begin{tcolorbox}[breakable,colback=teal!4,colframe=teal!35!black,title={Before you move on}]
What does \ta{கொட்டாவி} \ta{விடு} mean? (\textquotedblleft{}To yawn\textquotedblright{}.) Say it once more.
\end{tcolorbox}
